\section{Proofs of the exact results}
\label{app:proofs}
\paragraph{Setting.}
Fix a target with scores $x\in[0,100]^n$ and rival rows $r_1,\dots,r_m$, which stay fixed. The scenario cap on task $t$ gives the effective cap $c_t=\max(x_t,\text{scenario cap}_t)\le100$ and headroom $h_t=c_t-x_t\ge0$; admissible plans are $y\in\prod_t[x_t,c_t]$, so $L_1=\sum_t(y_t-x_t)$, $L_\infty=\max_t(y_t-x_t)$ and $K=|\{t:y_t>x_t\}|$. For a score rule $f$ let $\tau=\max_j f(r_j)$ be the best rival value; the target is first iff $f(y)\ge\tau+\varepsilon$ (a target with $f(x)\ge\tau+\varepsilon$, or one already uniquely first, costs 0 by convention, and the formulas below concern the other targets). Since $\varepsilon>0$ and all rules below are continuous in $y$, the set of first-placing plans is closed, so every minimum is attained. Write $(u)^+=\max(0,u)$.

\begin{proposition}[Mean]\label{prop:mean}
Let $g=\tau+\varepsilon-\bar x>0$. The target can become first iff $\sum_t h_t\ge n g$. Then the minimum $L_1$ is $n g$, the minimum $K$ is the least $k$ such that the $k$ largest headrooms sum to at least $n g$, and the minimum $L_\infty$ is the least $d$ with $\sum_t\min(d,h_t)\ge n g$.
\end{proposition}
\begin{proof}
$\bar y\ge\tau+\varepsilon$ iff $\sum_t(y_t-x_t)\ge n g$, and $\sum_t(y_t-x_t)=L_1$ because $y\ge x$. Hence $L_1\ge ng$ for every first-placing plan, and any plan with $y_t-x_t\in[0,h_t]$ and sum exactly $ng$ attains it; such a plan exists iff $\sum_t h_t\ge ng$, which is the feasibility condition. A plan touching a set $S$ of tasks raises the sum by at most $\sum_{t\in S}h_t$, which is at most the sum of the $|S|$ largest headrooms; conversely any amount in $[0,\sum_{t\in S}h_t]$ is attainable on $S$, so the least $|S|$ is the stated prefix length. A plan with all rises at most $d$ gains at most $\sum_t\min(d,h_t)$, and every amount up to that is attainable, giving the $L_\infty$ statement.
\end{proof}

\begin{proposition}[Minimum]\label{prop:min}
Let $b=\tau+\varepsilon$ with $\tau$ the best rival minimum. The target can become first iff $c_t\ge b$ for all $t$. Then $L_1=\sum_t(b-x_t)^+$, $K=|\{t:x_t<b\}|$ and $L_\infty=\max_t(b-x_t)^+$, each the minimum of its measure.
\end{proposition}
\begin{proof}
$\min_t y_t\ge b$ iff $y_t\ge b$ for every $t$. For $t$ with $x_t\ge b$ this holds with $y_t=x_t$ (and then $c_t\ge x_t\ge b$). For $t$ with $x_t<b$ we need $y_t\ge b$, possible iff $c_t\ge b$, and the cheapest such value is $y_t=b$. The three measures are separable over tasks (sum, count, maximum) and each coordinate is forced to at least $\max(x_t,b)$, so the stated values are lower bounds attained by $y_t=\max(x_t,b)$. One coordinate with $c_t<b$ blocks first place regardless of the other coordinates.
\end{proof}

\begin{proposition}[Odd median]\label{prop:median}
Let $n=2q-1$ and $b=\tau+\varepsilon$ with $\tau$ the best rival median. Let $a_t=(b-x_t)^+$ if $c_t\ge b$ and $a_t=\infty$ otherwise, and $A=|\{t:x_t\ge b\}|$. The target can become first iff at least $q$ coordinates have $c_t\ge b$. Then the minimum $L_1$ is the sum of the $q$ smallest $a_t$, the minimum $L_\infty$ is the $q$th smallest $a_t$, and the minimum $K$ is $(q-A)^+$.
\end{proposition}
\begin{proof}
For odd $n$ the median is the $q$th largest coordinate, so $\operatorname{med}(y)\ge b$ iff at least $q$ coordinates satisfy $y_t\ge b$. Such a coordinate $t$ costs at least $a_t$ (and is impossible if $a_t=\infty$) in each of the three measures: $y_t-x_t\ge(b-x_t)^+$. For any first-placing plan let $S$ be a set of $q$ coordinates with $y_t\ge b$. Then $L_1\ge\sum_{t\in S}a_t\ge$ the sum of the $q$ smallest $a_t$, and $L_\infty\ge\max_{t\in S}a_t\ge$ the $q$th smallest $a_t$; raising exactly the $q$ cheapest coordinates to $\max(x_t,b)$ attains both. For $K$, the $A$ coordinates already at least $b$ cost nothing and count toward $q$; each touched coordinate adds at most one new coordinate at least $b$, so $K\ge(q-A)^+$, attained by raising $(q-A)^+$ coordinates with $x_t<b\le c_t$ to $b$, which is possible whenever the plan is feasible.
\end{proof}

\begin{proposition}[Even median]\label{prop:even}
Let $n=2q$ and $b=\tau+\varepsilon$. A plan makes the target first iff there is a threshold $a\ge b$ such that at least $q$ coordinates are at least $a$ and at least $q+1$ are at least $2b-a$. Minimum $L_1$ and $L_\infty$ over plans are attained at a threshold $a$ in the finite set $\{b\}\cup\{x_t,c_t,2b-x_t,2b-c_t\}$, and the minimum $K$ is found by enumerating pairs of thresholds from the values $\{x_t,c_t\}$.
\end{proposition}
\begin{proof}
Order $y$ decreasingly as $z_1\ge\dots\ge z_n$; the median is $(z_q+z_{q+1})/2$, which is at least $b$ iff $z_q+z_{q+1}\ge2b$. Given such a plan take $a=z_q\ge b$; then $q$ coordinates are $\ge a$ and $z_{q+1}\ge2b-a$ gives $q+1$ coordinates $\ge2b-a$. Conversely such an $a$ gives $z_q\ge a$ and $z_{q+1}\ge2b-a$, whose sum is at least $2b$. For a fixed $a$ the cheapest plan lifts a set $Q$ of $q$ coordinates to $a$ and one further coordinate $e$ to $2b-a$ (costs $(a-x_t)^+$ and $(2b-a-x_e)^+$, impossible if the cap is lower). For a fixed assignment $(Q,e)$ the cost is piecewise linear in $a$ with breakpoints at $x_t$ and $2b-x_e$, over the interval of $a$ where all caps are respected, whose endpoints are $b$, some $c_t$ or some $2b-c_t$. A piecewise-linear function on an interval attains its minimum at a breakpoint or an endpoint, so enumerating the listed values and optimizing the assignment (a sort) for each yields the optimum; $L_\infty$ is the same with the maximum in place of the sum. For $K$, raising a touched coordinate to its cap never hurts, so a plan is a touched set with $y_t=c_t$ on it and $y_t=x_t$ elsewhere; its middle pair takes values in $\{x_t,c_t\}$, and for each threshold pair from these values the number of coordinates that must be touched is a simple count, so the minimum over pairs is the minimum $K$. These solvers are compared with exhaustive linear programs in \texttt{analysis/tests/test\_scorerules.py}.
\end{proof}

\begin{proposition}[Cap threshold]\label{prop:cap}
Let $f$ be a rule for which the set of first-placing plans in the ceiling box is closed (all score rules above, and the ordinal rules on their finite level sets). Let $d^\ast$ be the minimum $L_\infty$ over first-placing plans under the ceiling cap, attained by closedness. Under the cap ``at most $+h$ points per task'' the target can become first iff $d^\ast\le h$.
\end{proposition}
\begin{proof}
The gain-$h$ box has effective caps $c_t=\max(x_t,\min(x_t+h,100))=\min(x_t+h,100)$, so it equals $\{y:x\le y\le100,\ y_t-x_t\le h\}$, the set of ceiling-box plans with $L_\infty\le h$. A first-placing plan exists in it iff some first-placing ceiling plan has $L_\infty\le h$, iff $d^\ast\le h$. No monotonicity of $f$ is needed.
\end{proof}

\begin{proposition}[Finite levels for ordinal rules]\label{prop:levels}
The Borda total, the plurality count and the pairwise task-majority relations that define the Condorcet requirement depend on the plan $y$ only through the comparisons of each $y_t$ with the rival scores $s_{r,t}$ on that task (below, equal, above). Let $\mathcal L$ be the set of plans in which every $y_t$ is the original $x_t$, equals a rival score $s_{r,t}\ge x_t$, or equals $s_{r,t}+\delta$ for such a rival score ($\delta$ the margin step). Each cost measure is minimised over $\mathcal L$ by an integer linear program. The minimum over $\mathcal L$ is an upper bound on the infimum over real plans, and first place being unreachable on $\mathcal L$ does not prove it unreachable over the reals.
\end{proposition}
\begin{proof}
A comparison outcome is constant on each open interval between consecutive distinct values of $\{x_t\}\cup\{s_{r,t}\}$ and on each such value. Within an outcome class the cheapest coordinate value is the lower end of the class; for ``strictly above $s$'' that end is an infimum that is not attained, which is why the margin step $\delta>0$ is a convention that makes a strict pass attainable. Values below $x_t$ are inadmissible and rivals scoring below $x_t$ are already beaten on that task, so only rival scores at or above $x_t$ matter. Every plan in $\mathcal L$ is a real admissible plan, hence the minimum over $\mathcal L$ is at least the infimum over real plans. With one-hot variables $z_{t,\ell}$ selecting a level of each task, the Borda total, the plurality count and each pairwise task-majority indicator are linear in $z$ (a rival's loss is the sum of the target's gains against it), so the costs are minimised by integer linear programming. The dynamic program of Section~\ref{sec:methods} solves the relaxation that keeps only the constraint against the rival of highest original total; if its witness satisfies all rival constraints it is optimal for the full problem, and infeasibility of the relaxation proves infeasibility on $\mathcal L$.
\end{proof}

\begin{proposition}[Constructed counterexamples]\label{prop:counter}
The three rows of Table~\ref{tab:counter} are exact.
\end{proposition}
\begin{proof}
By Propositions~\ref{prop:mean}--\ref{prop:median} with $n=5$, $q=3$. Row 1 (target $(10,0,0,0,0)$, rival $(3,3,3,3,3)$, cap 100): the mean needs $5.005$ points, available on the first task (headroom 90), so $K=1$; the median threshold is $b=3.001$ and the first coordinate already exceeds it, so two more must be lifted ($K=2$, $L_1=6.002$); the minimum needs all four others at $b$ ($K=4$, $L_1=12.004$). Row 2 (zero target, rival $(9,9,1,1,1)$, gain 5): the mean needs $5\times4.201=21.005$ points at $5$ per task, so $K=5$; the median and minimum thresholds are $1.001$, so $K=3$ ($L_1=3.003$) and $K=5$ ($5.005$). Row 3 (target $(4,4,4,4,0)$, rival $(9,9,1,1,1)$, caps $(5.1,5.1,5.1,5.1,1.1)$): the mean needs $5.005$ points at $1.1$ per task, so $K=5$; the target's median 4 already exceeds the rival's 1 (cost 0); the minimum needs only the last task lifted to $1.001$ ($K=1$). By $K$ the order is mean $<$ median $<$ minimum in row 1 but median $<$ minimum $<$ mean in row 3, and the cheapest rule is the mean in row 1 and the median in rows 2 and 3, so the ordering reverses.
\end{proof}
